\chapter{Breuken: alle vier de bewerkingen}\label{ch:g8:fractions}

Dit hoofdstuk maakt de rekenkunde van de \omterm{def:g6:fractions:def}{breuken} af: optellen en aftrekken met
\emph{elke} \omterm{def:g4:fractions:def}{noemer}, vermenigvuldigen, en --- de nieuwkomer --- delen, dat een
vermomde vermenigvuldiging blijkt te zijn. De tekens uit \cref{ch:g8:negprod} gaan
mee op reis.

\section{Optellen met willekeurige noemers}

\begin{method}[Gemeenschappelijke noemer]\label{met:g8:fractions:add}
Om $\frac ab$ en $\frac cd$ op te tellen of af te trekken:
\begin{enumerate}
  \item zoek een gemeenschappelijk \omterm{def:g5:division:multiple}{veelvoud} van $b$ en $d$ (het \omterm{def:g2:mult:def}{product}
        $b \times d$ werkt altijd; een kleiner \omterm{def:g5:division:multiple}{veelvoud} spaart werk);
  \item herschrijf beide \omterm{def:g6:fractions:def}{breuken} met die \omterm{def:g4:fractions:def}{noemer};
  \item tel de \omterm{def:g4:fractions:def}{tellers} op of trek ze af; vereenvoudig.
\end{enumerate}
\end{method}

\begin{example}\label{ex:g8:fractions:add}
Reken $\dfrac{5}{6} + \dfrac{3}{4}$ uit. Een gemeenschappelijk \omterm{def:g5:division:multiple}{veelvoud} van $6$ en
$4$ is $12$:
\[
\frac{5}{6} = \frac{10}{12},
\qquad
\frac{3}{4} = \frac{9}{12},
\qquad
\frac{5}{6} + \frac{3}{4} = \frac{10 + 9}{12} = \frac{19}{12}.
\]
Met tekens: $\dfrac{2}{3} - \dfrac{7}{5} = \dfrac{10}{15} - \dfrac{21}{15} =
-\dfrac{11}{15}$.
\end{example}

\section{Breuken vermenigvuldigen}

\begin{theorem}[Product van breuken]\label{thm:g8:fractions:mult}
\[
\frac{a}{b} \times \frac{c}{d} = \frac{a \times c}{b \times d} .
\]
\end{theorem}

\begin{proof}[Waarom, op een plaatje]
Neem $\frac34$ van $\frac25$ van een vierkant: knip het vierkant in $5$ verticale
stroken en houd er $2$; knip horizontaal in $4$ en houd $3$ rijen van wat er
overbleef. Wat je overhoudt, is een rooster van $3 \times 2$ vakjes van de
$4 \times 5$: de \omterm{def:g6:fractions:def}{breuk} $\frac{6}{20}$.
\end{proof}

\begin{omfigure}
\begin{tikzpicture}[scale=0.68]
  \draw[gray!50, thin] (0,0) grid[xstep=1, ystep=1] (5,4);
  \fill[omDef!18] (0,0) rectangle (2,4);
  \fill[omThm!35] (0,0) rectangle (2,3);
  \draw[thick] (0,0) rectangle (5,4);
  \draw[omDef, thick] (2,0) -- (2,4);
  \draw[omThm, thick] (0,3) -- (5,3);
  \node[below, font=\small, omDef] at (1,0) {$\frac25$ gehouden};
  \node[right, font=\small, omThm] at (5.05,1.5) {$\frac34$ daarvan};
\end{tikzpicture}

{\small $\frac34 \times \frac25$: het dubbel gekleurde gebied is
$3 \times 2 = 6$ vakjes van de $4 \times 5 = 20$, dus
$\frac{6}{20} = \frac{3}{10}$ --- de \omterm{def:g4:fractions:def}{tellers} vermenigvuldigd, de \omterm{def:g4:fractions:def}{noemers}
vermenigvuldigd.}
\end{omfigure}

\begin{example}\label{ex:g8:fractions:mult}
Vereenvoudig \emph{voordat} je vermenigvuldigt, door gemeenschappelijke factoren weg
te strepen:
\[
\frac{7}{12} \times \frac{9}{14}
= \frac{7 \times 9}{12 \times 14}
= \frac{7 \times 3 \times 3}{3 \times 4 \times 2 \times 7}
= \frac{3}{8}
\]
(de factor $7$ en één factor $3$ vallen boven en onder weg). Met tekens gelden eerst
de regels van \cref{thm:g8:negprod:rules}:
$\left(-\frac23\right) \times \left(-\frac{5}{4}\right) = +\frac{10}{12} =
\frac56$.
\end{example}

\section{Breuken delen}

\begin{definition}[Omgekeerde]\label{def:g8:fractions:inverse}
De \emph{omgekeerde}\index{omgekeerde} van een getal $x$ dat niet \omterm{def:g1:counting:zero}{nul} is, is het getal
dat met $x$ vermenigvuldigd $1$ geeft. De omgekeerde van $\frac ab$ is $\frac ba$,
want $\frac ab \times \frac ba = \frac{ab}{ab} = 1$; de omgekeerde van een geheel
getal $n$ is $\frac1n$.
\end{definition}

\begin{theorem}[Delen is vermenigvuldigen met de omgekeerde]\label{thm:g8:fractions:div}
Voor $c \neq 0$:
\[
\frac{a}{b} \div \frac{c}{d} = \frac{a}{b} \times \frac{d}{c} .
\]
\end{theorem}

\begin{proof}
Per definitie is het \omterm{def:g3:division:remainder}{quotiënt} $q$ van $\frac ab$ door $\frac cd$ het getal waarvoor
$q \times \frac cd = \frac ab$. Test de kandidaat $q = \frac ab \times \frac dc$:
\[
\left(\frac ab \times \frac dc\right) \times \frac cd
= \frac ab \times \left(\frac dc \times \frac cd\right)
= \frac ab \times 1 = \frac ab . \qedhere
\]
\end{proof}

\begin{example}\label{ex:g8:fractions:div}
\[
\frac{3}{5} \div \frac{7}{10}
= \frac{3}{5} \times \frac{10}{7}
= \frac{3 \times 10}{5 \times 7}
= \frac{30}{35} = \frac{6}{7}.
\]
Een controle op de grootte: delen door $\frac{7}{10}$, een getal kleiner dan $1$, moet
een uitkomst geven die \emph{groter} is dan $\frac35$ --- en $\frac67 > \frac35$ is
inderdaad zo ($\frac{30}{35} > \frac{21}{35}$).
\end{example}

\begin{example}[Breukstrepen in breukstrepen]\label{ex:g8:fractions:complex}
Een ``dubbeldekker'' is niets anders dan een deling die verticaal geschreven is:
\[
\frac{\ \frac{2}{3}\ }{\ \frac{5}{6}\ }
= \frac{2}{3} \div \frac{5}{6}
= \frac{2}{3} \times \frac{6}{5}
= \frac{12}{15} = \frac{4}{5}.
\]
Zoek eerst de \emph{hoofdstreep} (de langste) en pas dan
\cref{thm:g8:fractions:div} toe.
\end{example}

\begin{method}[Gemengde berekeningen]\label{met:g8:fractions:mixed}
In een uitdrukking met \omterm{def:g6:fractions:def}{breuken} en de vier bewerkingen:
\begin{enumerate}
  \item pas de gewone voorrangsregels toe (\cref{ch:g7:priorities});
  \item zet elke deling om in een vermenigvuldiging met de \omterm{def:g8:fractions:inverse}{omgekeerde};
  \item beslis eerst de tekens (\cref{met:g8:negprod:compute}) en vermenigvuldig
        daarna, of zoek gemeenschappelijke \omterm{def:g4:fractions:def}{noemers};
  \item vereenvoudig bij de eerste gelegenheid.
\end{enumerate}
\end{method}

\begin{example}\label{ex:g8:fractions:mixed}
\begin{align*}
\frac12 + \frac{2}{3} \times \frac{9}{4}
&= \frac12 + \frac{2 \times 9}{3 \times 4}
&& \text{(eerst het product!)} \\
&= \frac12 + \frac{18}{12} = \frac12 + \frac32
= \frac{4}{2} = 2 .
\end{align*}
\end{example}

\section{Oefeningen}

\begin{exercise}[$\star$]\label{exo:g8:fractions:1}
Reken uit en vereenvoudig:
\[
\frac{3}{4} + \frac{2}{5}, \qquad
\frac{5}{6} - \frac{3}{8}, \qquad
\frac{7}{10} + \frac{8}{15} .
\]
\end{exercise}

\begin{exercise}[$\star$]\label{exo:g8:fractions:2}
Reken uit (eerst de tekens):
\[
-\frac{2}{7} + \frac{3}{14}, \qquad
\frac{1}{6} - \frac{5}{4}, \qquad
-\frac{3}{5} - \frac{1}{2} .
\]
\end{exercise}

\begin{exercise}[$\star$]\label{exo:g8:fractions:3}
Reken uit, en vereenvoudig voordat je vermenigvuldigt:
\[
\frac{5}{8} \times \frac{4}{15}, \qquad
\frac{9}{14} \times \frac{7}{3}, \qquad
\left(-\frac{6}{5}\right) \times \frac{10}{9} .
\]
\end{exercise}

\begin{exercise}[$\star$]\label{exo:g8:fractions:4}
Geef de \omterm{def:g8:fractions:inverse}{omgekeerde} van: $\dfrac{3}{7}$; \quad $5$; \quad $\dfrac{1}{4}$; \quad
$-\dfrac{2}{9}$.
\end{exercise}

\begin{exercise}[$\star$]\label{exo:g8:fractions:5}
Reken uit:
\[
\frac{4}{9} \div \frac{2}{3}, \qquad
\frac{7}{5} \div 14, \qquad
6 \div \frac{3}{4} .
\]
\end{exercise}

\begin{exercise}[$\star$]\label{exo:g8:fractions:6}
Reken de dubbeldekkers uit:
\[
\frac{\ \frac{3}{4}\ }{\ \frac{9}{8}\ },
\qquad
\frac{\ \frac{5}{6}\ }{\ 10\ },
\qquad
\frac{\ 2\ }{\ \frac{4}{7}\ } .
\]
\end{exercise}

\begin{exercise}[$\star$]\label{exo:g8:fractions:7}
Reken uit, met respect voor de voorrangsregels:
\[
\frac{1}{3} + \frac{1}{2} \times \frac{4}{5},
\qquad
\left(\frac{1}{3} + \frac{1}{2}\right) \times \frac{4}{5} .
\]
\end{exercise}

\begin{exercise}[$\star\star$]\label{exo:g8:fractions:8}
Een tank is voor $\frac{2}{5}$ gevuld. Er wordt $\frac{1}{3}$ \emph{van de \omterm{def:g2:measure:units}{inhoud} van
de tank} bij gedaan. Welke \omterm{def:g6:fractions:def}{breuk} van de tank is nu gevuld? En welke \omterm{def:g6:fractions:def}{breuk} is nog
leeg?
\end{exercise}

\begin{exercise}[$\star\star$]\label{exo:g8:fractions:9}
Driekwart van de leerlingen van een klas haalde een toets; van hen kreeg twee derde
meer dan $14$ van de $20$. Welke \omterm{def:g6:fractions:def}{breuk} van de klas kreeg meer dan $14$? Als dat $12$
leerlingen zijn, hoe groot is de klas dan?
\end{exercise}

\begin{exercise}[$\star\star$]\label{exo:g8:fractions:10}
Een touw van $\frac{15}{2}$ m moet in stukken van elk $\frac{3}{4}$ m geknipt worden.
Hoeveel stukken krijg je? (Een deling van \omterm{def:g6:fractions:def}{breuken} --- ga na dat het antwoord een heel
getal is.)
\end{exercise}

\begin{exercise}[$\star\star$]\label{exo:g8:fractions:11}
Reken uit:
\[
E = \frac{2}{3} - \frac{5}{6} \div \frac{10}{9},
\qquad
F = \left(-\frac{1}{2}\right)^2 \times \frac{8}{3} .
\]
\end{exercise}

\begin{exercise}[$\star\star\star$]\label{exo:g8:fractions:12}
Vereenvoudig de uitdrukking
\[
1 + \cfrac{1}{1 + \cfrac{1}{1 + 1}}
\]
(begin bij de binnenste \omterm{def:g6:fractions:def}{breuk} en werk naar buiten, één streep per keer).
\end{exercise}

\section{Opgave: Egyptische breuken}

\begin{problem}[{Weekendopgave --- elke breuk is een som van verschillende
stambreuken}]
\label{pb:g8:fractions:1}
Een \emph{stambreuk} is een \omterm{def:g6:fractions:def}{breuk} met \omterm{def:g4:fractions:def}{teller} $1$. De schrijvers van het oude Egypte ---
de Rhind-papyrus werd rond $1550$~v.Chr. gekopieerd --- schreven elke \omterm{def:g6:fractions:def}{breuk} als een \omterm{def:g1:addition:def}{som}
van \emph{verschillende} stambreuken: nooit $\frac27 = \frac17 + \frac17$, maar
$\frac27 = \frac14 + \frac{1}{28}$. Zo'n \omterm{def:g1:addition:def}{som} heet een \emph{Egyptische schrijfwijze}
van de \omterm{def:g6:fractions:def}{breuk}. In deze opgave werk je de optel- en vergelijktechnieken van dit
hoofdstuk (\cref{met:g8:fractions:add}) uit tot een methode --- in $1202$ door
Fibonacci gepubliceerd --- die van elke \omterm{def:g6:fractions:def}{breuk} tussen $0$ en $1$ een Egyptische
schrijfwijze oplevert, en eindig je met een gevierde erfenispuzzel.

\medskip
\noindent\textbf{Deel I --- Stambreuken en de splitsingsformule.}
\begin{enumerate}
  \item Reken $\frac12 + \frac13$ en $\frac12 + \frac14$ uit, en leid Egyptische
        schrijfwijzen van $\frac56$ en $\frac34$ af.
  \item Ga de vermelding van de schrijvers voor $\frac27$ na: laat zien dat
        $\frac14 + \frac{1}{28} = \frac27$.
  \item Ga na dat $\frac13 + \frac16 = \frac12$ en dat
        $\frac14 + \frac{1}{12} = \frac13$. Bewijs daarna de
        \emph{splitsingsformule}: voor elk heel getal $n \geq 1$ geldt
        \[
        \frac{1}{n} = \frac{1}{n+1} + \frac{1}{n(n+1)} .
        \]
  \item Leid een Egyptische schrijfwijze van $1$ zelf af, als \omterm{def:g1:addition:def}{som} van drie
        verschillende stambreuken.
  \item Splits de laatste term van $\frac56 = \frac12 + \frac13$ om een \emph{tweede}
        Egyptische schrijfwijze van $\frac56$ te krijgen, met drie stambreuken. Leg uit
        waarom geen enkele \omterm{def:g6:fractions:def}{breuk} een unieke Egyptische schrijfwijze heeft.
\end{enumerate}

\medskip
\noindent\textbf{Deel II --- De hebberige methode van Fibonacci.}
De methode is: trek de \emph{grootste} stambreuk af die past, en herhaal met wat
overblijft.
\begin{enumerate}[resume]
  \item Laat met gemeenschappelijke \omterm{def:g4:fractions:def}{noemers} zien dat $\frac14 < \frac27 < \frac13$.
        Waarom bewijst dat dat $\frac14$ de grootste stambreuk kleiner dan $\frac27$
        is?
  \item Reken $\frac27 - \frac14$ uit. Welke schrijfwijze van $\frac27$ levert de
        hebberige methode op?
  \item Neem nu $\frac57$. Laat zien dat de grootste stambreuk kleiner dan $\frac57$
        gelijk is aan $\frac12$, en reken $\frac57 - \frac12$ uit.
  \item Laat zien dat de grootste stambreuk kleiner dan $\frac{3}{14}$ gelijk is aan
        $\frac15$ (vergelijk $\frac{3}{14}$ met $\frac14$ en met $\frac15$), reken
        $\frac{3}{14} - \frac15$ uit, en besluit:
        \[
        \frac57 = \frac12 + \frac15 + \frac{1}{70} .
        \]
        Ga die schrijfwijze rechtstreeks na met de gemeenschappelijke \omterm{def:g4:fractions:def}{noemer} $70$.
  \item Laat de hebberige methode op $\frac45$ lopen, ga bij elke stap na dat je
        stambreuk de grootste is die past, en controleer de uiteindelijke schrijfwijze.
\end{enumerate}

\medskip
\noindent\textbf{Deel III --- Waarom de methode altijd stopt, en de achttiende
kameel.}
\begin{enumerate}[resume]
  \item In de vragen 7--10 is de methode toegepast op $\frac27$, $\frac57$ en
        $\frac45$. Noem voor elk de \omterm{def:g4:fractions:def}{tellers} van de opeenvolgende \omterm{def:g3:division:remainder}{resten} (voordat je
        vereenvoudigt). Wat valt je op?
  \item Laat voor een \omterm{def:g6:fractions:def}{breuk} $\frac ab$ en een heel getal $n$ zien dat
        \[
        \frac{a}{b} - \frac{1}{n}
        = \frac{a \times n - b}{b \times n} .
        \]
  \item Stel dat de hebberige methode $\frac1n$ van $\frac ab$ aftrekt: dan past
        $\frac1n$, maar is $\frac{1}{n-1}$ al te groot, dus
        $\frac{1}{n-1} > \frac ab$. Zet beide op de gemeenschappelijke \omterm{def:g4:fractions:def}{noemer}
        $b \times (n - 1)$ en leid af dat $b > a \times n - a$, en dus dat de nieuwe
        \omterm{def:g4:fractions:def}{teller} voldoet aan
        \[
        a \times n - b < a .
        \]
        Leg uit waarom dat bewijst dat de methode altijd stopt.
  \item Een oud verhaal: een vader sterft en laat $17$ kamelen achter, met in zijn
        testament $\frac12$ van de kudde voor zijn oudste kind, $\frac13$ voor het
        tweede en $\frac19$ voor het derde --- en geen van de delen is een heel aantal
        kamelen. Een buurman leent het gezin een achttiende kameel; de kinderen nemen
        $9$, $6$ en $2$ kamelen, en de geleende kameel gaat terug. Reken
        $\frac12 + \frac13 + \frac19$ uit en verklaar de truc: waarom kon de kudde in
        hele kamelen verdeeld worden, en waarom kreeg geen enkel kind minder dan het
        testament beloofde?
  \item Repareer het testament: stel drie \emph{verschillende} stambreuken voor die
        samen precies $1$ zijn, en geef de kleinste kudde waarvoor alle drie de delen
        hele aantallen kamelen zijn.
\end{enumerate}
\end{problem}
